\chapter{Машинное обучение: первые модели}
\label{ch:ml}

\textit{[В этой главе планируется разобрать:
линейную регрессию, логистическую регрессию и метод опорных векторов.
Каждая модель сводится к задаче оптимизации, которую можно решать
методами, изученными в главе~\ref{ch:opt}, в~частности — методом Ньютона.]}

\section{Линейная регрессия}
\textit{[Заготовка раздела.]}

\section{Логистическая регрессия и её обучение методом Ньютона}
\textit{[Заготовка: применение теоремы~\ref{thm:newton-rate} к функции
правдоподобия логистической регрессии — она сильно выпуклая, поэтому
метод Ньютона сходится квадратично. Эта схема исторически известна
как «итерационно перевзвешенные наименьшие квадраты», IRLS.]}

\section{Глубокие нейронные сети: краткий очерк}
\textit{[Заготовка: backpropagation, оптимизатор Adam, упоминание
оптимизатора Muon из раздела~\ref{sec:newton}.]}
